\section{Restriction to zero or the newest parameter}
\label{sec:newest-witness}

Let $\mathcal R_{\mathrm{new}}$ be the restricted calculus of
Section~\ref{sec:restricted-witness}, with variable witnesses further
restricted to the newest numerical parameter of the current scope.
The literal $0$ remains available.  At an empty scope only $0$ is
available.  Introduction opens its step variable \emph{after} selecting
the witness, so the newest parameter here means the newest parameter
before that step variable is opened.

This restriction is not invariant under arbitrary weakening of a
derivation by a new numerical parameter.  Accordingly, the direct
structural proof from the preceding section cannot be reused.  The
following argument uses a protected translation instead.  It assumes
that the predecessor identity has been proved in the equational
computation/progress fragment, whose derivations are invariant under
numerical renaming; this is stronger than merely assuming an arbitrary
derivation of that identity in the target calculus.

\subsection{Hereditary derivability and uniform cut}

Write $\Gamma\Vdash A$ if, for every injective renaming $\rho$ of the
numerical parameters into any larger numerical scope and every
formula context $\Delta$ containing $\rho\Gamma$, there is an
$\mathcal R_{\mathrm{new}}$ derivation of $\rho A$ from $\Delta$.
This is a metatheoretic family of actual target derivations, not an
additional rule or a second object language.  It is closed by definition
under injective renaming and formula-context weakening.  Membership
$A\in\Gamma$ implies $\Gamma\Vdash A$.

\begin{lemma}[Uniform cut]
If $\Pi$ is an actual target derivation of $A$ from $\Gamma$, and
$\Delta\Vdash B$ for every $B\in\Gamma$ in the same numerical scope,
then there is an actual target derivation of $A$ from $\Delta$.
\end{lemma}
\begin{proof}
Induct on $\Pi$.  At a hypothesis use its given family at the identity
renaming.  At an eigenvariable rule, each replacement assumption has
a new derivation in the enlarged numerical scope because its proof is
hereditary.  Add the temporary assumption as a hypothesis and apply the
induction hypothesis to the subderivation.  The numerical scope and the
chosen witness of the original target rule have not changed, so a
newest-variable witness is still admissible.  No existing target proof
is renamed under the new eigenvariable.  Contexts are scoped: $\Delta$
has the same numerical scope as $\Gamma$, so it cannot already mention
the new eigenparameter of $\Pi$.  Thus its freshness condition is
preserved as well.
\end{proof}

As a consequence, cut is admissible for $\Vdash$: if
$\Gamma\Vdash A$ and $\Gamma,A\Vdash B$, then $\Gamma\Vdash B$.
At each renamed extension, apply uniform cut to the second derivation
using the hereditary family for the first.  More generally a finite
list of hereditary replacement assumptions can be cut simultaneously.

\subsection{The logical and equational interfaces at the hereditary level}

Computation, crossed transport, and QE preserve $\Vdash$.  They introduce
no new numerical parameter, and their conclusion is obtained by applying
the same primitive rule to the available members of the proof families.
Native DE also preserves $\Vdash$: at each renamed extension open the
witness at the newest fresh parameter, and instantiate the hereditary
family for its second premise at the correspondingly extended renaming.
The original eigenparameter maps to this newly chosen parameter.

The closure rules QC and QP hold at the hereditary level.  For QC,
expand the proof through the neutral theorem $D_{r,w}(0=0)$ as in the
preceding section.  Its fresh witness $b$ is newest, and hence is
admissible for the dummy-witness case introduction.  The given step
family is instantiated in the numerical scope containing $b$ and a
new step parameter $z$, rather than obtained by renaming a single
target derivation.  QP uses the literal witness $0$ and is unchanged.

Full formula transport holds for hereditary proofs as well.  The
induction on the formula is exactly the explicit induction in the
preceding section, with hereditary cut in place of ordinary cut.
In the $D$ case the witness that is reintroduced is the one just
opened, so it is newest.  In the $Q$ case use hereditary QC.
All equality and closure premises required below another binder are
regenerated from their hereditary families.  Thus
\[
 \Gamma\Vdash s=t,\quad\Gamma\Vdash H(s)
 \quad\Longrightarrow\quad\Gamma\Vdash H(t).
 \tag{HL}
\]
This statement is uniform in all actual term assignments to free
slots of $H$, as in the previous transport proof.

\subsection{A hereditary naming packet}

Retain $T(v)=D_{x,y}(C(y,v,C(v,0,Sx))=v)$.
The theorem $T(0)$ has a derivation in every numerical scope by
progressive introduction at $0$.
There is also a derivation of $Q_vT(v)$ in every numerical scope.
To see this, \emph{expand} QC instead of invoking it on a supposedly
hereditary proof of $T(Sz)$: open the neutral theorem at a fresh $b$;
introduce the dummy-witness $D$ at $b$; and in its step premise, with
$z$ newly opened, prove $T(Sz)$ by case introduction at $z$.
Here $z$ really is newest, even though $b$ has already been added.
The successor row is a computation theorem.  Eliminate the two dummy
witnesses exactly as before.

This construction works afresh in every ambient numerical scope and
formula context.  Consequently $\varnothing\Vdash Q_vT(v)$ and
$\varnothing\Vdash T(0)$.  QE at an arbitrary actual term $t$ then
gives
\[
 \varnothing\Vdash T(t).
\]
No claim that $T(Sz)$ has a hereditary direct-introduction proof was
needed; such a claim would incorrectly allow $z$ to cease being newest.

Opening $T(St)$ at a fresh newest $b$ gives $b=t$ by computations and
the predecessor identity.  In the opened context, that equality has
a hereditary proof: its derivation uses only the retained closure,
QE, and equational derivations in the renaming-stable fragment.
Therefore it can replace an assumed $b=t$ in an arbitrary target
subderivation by uniform cut.  The packet's DE then discharges $b$.

\subsection{Translation of the native kernel}

\begin{theorem}
The native scalar kernel and $\mathcal R_{\mathrm{new}}$ derive
exactly the same sequents.  Thus in the combined witness-mode rule it
is sufficient to permit the literal $0$ and the newest numerical
parameter, and no other witness.
\end{theorem}
\begin{proof}
Induct on a native-kernel derivation to prove the stronger assertion
$\Gamma\Vdash A$.  All cases except native case introduction follow
from the hereditary interfaces just established.

Consider a case introduction at an arbitrary term $t$, whose premise
is $\Gamma,B(t,0)\vdash B(t,Sz)$.  Its induction hypothesis is a
hereditary family.  To prove the conclusion at a chosen renaming
$\rho$ and context extension $\Delta$, open the hereditary packet
$T(S\rho t)$ at the newest fresh parameter $b$.  It supplies the
hereditarily derivable equation $b=\rho t$.

Temporarily retain $b=\rho t$ as an assumption in addition to
$\Delta$ weakened to the scope containing $b$.  Construct a target
case introduction of the desired $D$ formula at this newest witness
$b$.  In its premise, under the temporary assumption $B(b,0)$ and a
newly opened step parameter, hereditary
transport gives $B(\rho t,0)$.  Instantiate the source premise's
hereditary family at the injection sending the source step parameter
to this new step parameter.  Hereditary cut supplies its missing
zero-row assumption, yielding $B(\rho t,Sz)$.  Hereditary transport
backwards gives $B(b,Sz)$.  Take the member of this step family at the
identity renaming and apply the primitive target introduction at $b$.
Finally use uniform cut to replace the retained assumption
$b=\rho t$ by its hereditary proof from the opened packet; the other
assumptions are supplied by membership in the opened context.  Then
discharge the packet.  The resulting formula is independent of
$b$.

The construction was performed for every $\rho,\Delta$, so it proves
the required hereditary conclusion.  The reverse implication is
immediate: every target primitive rule is already a derived rule of
the native kernel.  Taking the identity member of the hereditary
family completes the comparison.
\end{proof}

The result removes the need to supply an arbitrary term or a variable
index as an introduction witness.  A checker need only choose between
two fixed witnesses, $0$ and the newest parameter, rejecting the
second choice at an empty scope.  This is a statement about the
available witness data, not by itself a claim that a particular
serialization or complete checker is smaller.

\subsection{A total two-choice interface}

For an implementation it is simpler to make both choices total.  Let
$\beta\in\{0,1\}$ be the only introduction datum, and in a numerical
scope of size $n$ put
\[
 t_{n,\beta}:=
 \begin{cases}
 0,&\beta=0\text{ or }n=0,\\
 b_{\mathrm{new}},&\beta=1\text{ and }n>0,
 \end{cases}
 \qquad
 \eta_{B,\beta}:=
 \begin{cases}
 z,&B\text{ is an equation and }\beta=0,\\
 0,&\text{otherwise}.
 \end{cases}
\]
The introduction is the single schema
\[
 \frac{\Gamma,B(t_{n,\beta},\eta_{B,\beta})
        \vdash B(t_{n,\beta},Sz)}
      {\Gamma\vdash D_{x,y}B(x,y)}.
 \tag{Bit-DI}
\]
The mode is selected by the bit, not by inspecting the chosen term.
Thus the choice $\beta=1$ in an empty scope is case introduction at
$0$, including when the matrix is an equation.

\begin{corollary}[One-bit introduction]
Replacing RDI by Bit-DI gives exactly the native scalar kernel's
derivable sequents.
\end{corollary}
\begin{proof}
Every rule of $\mathcal R_{\mathrm{new}}$ is an instance of Bit-DI.
Every instance of Bit-DI is either ordinary case introduction or
equational progress at witness $0$, hence is derived in the native
kernel.  The comparison theorem sandwiches the new calculus between
two calculi with the same derivability relation.
\end{proof}

At the certificate level one may require a literal Boolean datum, or
define a total projection of an arbitrary JSON value onto a bit.
Either is a serialization choice that must be stated explicitly.
The mathematical certificate field no longer represents a term or a
variable index.  In particular its validation never needs to traverse
a proposed witness term.  Replacing the empty-scope rejection by the
specified zero fallback is sound by the corollary; it is not an
unchecked exception to the syntax.
